Activation Addition (ActAdd) enables targeted control of large language model (LLM) behavior at inference time by intervening on internal activations rather than model weights \citep{turnerSteeringLanguageModels2023}.
Contrastive Activation Addition (CAA) provides a contrastive-pair variant of the same intervention approach \citep{rimskySteeringLlama22024}.
These methods are computationally inexpensive relative to fine-tuning, yet output quality can degrade under strong steering in ways that prior standard metrics, such as perplexity, do not capture.

According to \cite{tanAnalysingGeneralisationReliability2024}, steerability varies across inputs and prompt settings (such as different system prompts).
Likewise, \cite{vonrutteLanguageModelsGuide2024} shows that strong concept detection does not necessarily provide better guidance.
Conditional activation steering further shows that behavior vectors can be applied selectively by context rather than uniformly \citep{leeProgrammingRefusalConditional2025}.
Despite these advances, the representation engineering literature still lacks a systematic taxonomy of qualitative side effects at scale.

When deploying these models in real-world settings, understanding these unmeasured side effects is necessary, as they present practical problems that extend beyond standard benchmark scores.
A model can remain fluent while reversing factual answers, inflating disclaimers, or shifting target-token logits without changing the generated text. These are instances of side effects that perplexity and benchmark performance scores cannot flag.
Without a clear framework to track these side effects, adjusting the steering multiplier essentially becomes a guessing game. If the multiplier ($|\alpha|$) is pushed too high, it introduces unwanted side effects, but if it is too low, the target behavior never actually appears.

This thesis inquires into which structural side effects occur under activation steering across nine multipliers and how severe they are.
We treat side effects as a reproducible measurement problem: 19 dimensions scored on an ordinal 0-5 severity scale relative to an $\alpha{=}0$ baseline.
The full taxonomy appears in Table~\ref{tab:side_effect_taxonomy}.
We apply CAA to Llama-3.1-8B-Instruct across 36 Model-Written Evaluations datasets \citep{perezDiscoveringLanguageModel2022}, sweep $\alpha \in \{-2.0,\ldots,+2.0\}$, and evaluate outputs with a two-pass LLM-as-a-judge pipeline: Pass~1 screens for steerability and broad anomaly categories; Pass~2 scores severity on flagged samples across the 19 side-effect dimensions.

This thesis makes three contributions.
First, we introduce a taxonomy of steering side effects (including steering asymmetry, token inflation, factual reversal, and inverse logit polarity) with ordinal scoring rubrics and a two-pass annotation protocol.
Second, we scale this evaluation to 28{,}807 samples and distinguish primary claims from exploratory findings using pre-specified practical-significance criteria (Section~\ref{sec:Statistical Analysis}).
Table~\ref{tab:dimension_claims} summarizes the resulting primary vs.\ exploratory classifications.
Third, we validate the measurement protocol with inter-rater agreement on a 1{,}800-sample subsample, per-$\alpha$ severity checks, human rating on the four primary dimensions, and exploratory case studies of logit-text decoupling.

Qualitative case studies and inter-rater agreement details appear in the appendices.
